%% Online appendix for frl_three_results.tex
%% Compile with: latexmk -pdf frl_three_results_online_appendix.tex

\documentclass[12pt]{article}

\usepackage[margin=1in]{geometry}
\usepackage{amsmath}
\usepackage{booktabs}
\usepackage{graphicx}
\usepackage{threeparttable}
\usepackage{longtable}
\usepackage{array}
\usepackage{float}
\usepackage{xurl}
\graphicspath{{figures/}}

\newcommand{\pp}{\,pp}
\newcommand{\sym}[1]{\ifmmode^{#1}\else\(^{#1}\)\fi}
\renewcommand{\thetable}{A\arabic{table}}
\renewcommand{\thefigure}{A\arabic{figure}}

\title{Online Appendix to\\
\textit{Who Gains from Generative AI Model Releases?}}
\author{Zhuo Chen \and Ruishuang Lu}
\date{}

\begin{document}
\maketitle

\section{Sample Construction and Position Coding}

The event universe combines public model chronologies with Artificial
Analysis benchmark records. Same-day releases from the same model family
are consolidated. Announcement dates are verified against vendor blogs,
release notes, repositories, page metadata, and archived first captures.
The pipeline retains 125 day-verified events and excludes one event whose
identity is inconsistent across the two source systems. The return analysis
uses 124 events. The annual counts are 3, 8, 45, 59, and 9 from 2022 through
2026.

The firm universe is the May 1, 2026 constituent list of the Nasdaq-100
Technology Sector Index. It contains 45 securities and 44 issuers. GOOG and
GOOGL are retained as distinct traded securities and are combined at the
issuer level in a robustness check. The balanced event-firm panel has 5,625
rows before the identity exclusion and regression-data requirements. The
main position regression contains 5,441 observations, 124 events, and 45
securities.

Two coders independently evaluate 45 securities against 25 model
developers. The resulting 1,125 firm--developer pairs contain eight binary
position indicators. Table~\ref{tab:a1} reports agreement before
adjudication. Overall agreement is 98.5\% across 9,000 binary cells.
Disagreements are resolved under the written codebook and retained in the
decision files.

\begin{table}[H]
\centering
\begin{threeparttable}
\caption{Independent coding agreement}\label{tab:a1}
\small
\begin{tabular}{lrrrr}
\toprule
Position & Agreement (\%) & Cohen's $\kappa$ &
Disagreements & Cells \\
\midrule
Hardware supplier      & 100.0 & 1.000 & 0  & 1,125 \\
Cloud platform         & 100.0 & 1.000 & 0  & 1,125 \\
Downstream integrator  & 95.6  & 0.902 & 49 & 1,125 \\
Downstream deployer    & 97.6  & 0.835 & 27 & 1,125 \\
Downstream enabler     & 98.2  & 0.328 & 20 & 1,125 \\
Competitor             & 96.7  & 0.823 & 37 & 1,125 \\
Investor               & 100.0 & 1.000 & 0  & 1,125 \\
Listed owner           & 100.0 & 1.000 & 0  & 1,125 \\
\midrule
All binary cells       & 98.5  &       & 133 & 9,000 \\
\bottomrule
\end{tabular}
\begin{tablenotes}\footnotesize
\item Notes: Agreement is computed on the 45-security universe before
adjudication. The enabler indicator is sparse, so its high raw agreement
coexists with a lower prevalence-adjusted $\kappa$.
\end{tablenotes}
\end{threeparttable}
\end{table}

The capability sample contains 87 language-model events with comparable
Artificial Analysis intelligence scores. The index has a standard
deviation of 12.7 points. The remaining 37 events lack a comparable
intelligence score. A separate media-generation flag identifies 41 events;
four multi-component releases contain both a scored language model and a
media model.

\section{Position Results Across Benchmarks and Fixed Effects}

Table~\ref{tab:a2} reports the two positive positions across the full set of
return benchmarks. QQQ removes broad Nasdaq returns. SOXX absorbs the
performance of the semiconductor sector. Event fixed effects compare firms
inside the same release. Hardware and competitor coefficients remain close
to their baseline magnitudes in each specification.

\begin{table}[H]
\centering
\begin{threeparttable}
\caption{Hardware and competitor coefficients across return models}
\label{tab:a2}
\small
\setlength{\tabcolsep}{4pt}
\begin{tabular}{llrrr}
\toprule
Specification & Position & Coef. (pp) & Event CR2 $p$ &
Event $\times$ firm $p$ \\
\midrule
S\&P 500 market model
 & Hardware   & 3.476\sym{***} & 0.0017 & 0.0145 \\
 & Competitor & 3.385\sym{***} & 0.0003 & 0.0222 \\
\addlinespace
Fama--French three-factor
 & Hardware   & 3.428\sym{***} & 0.0021 & 0.0159 \\
 & Competitor & 3.371\sym{***} & 0.0003 & 0.0368 \\
\addlinespace
QQQ market model
 & Hardware   & 3.459\sym{***} & 0.0016 & 0.0143 \\
 & Competitor & 3.372\sym{***} & 0.0003 & 0.0253 \\
\addlinespace
SOXX market model
 & Hardware   & 2.898\sym{***} & 0.0027 & 0.0200 \\
 & Competitor & 3.182\sym{***} & 0.0007 & 0.0282 \\
\addlinespace
S\&P 500 with event fixed effects
 & Hardware   & 3.364\sym{***} & 0.0023 & 0.0282 \\
 & Competitor & 3.259\sym{***} & 0.0004 & 0.0408 \\
\bottomrule
\end{tabular}
\begin{tablenotes}\footnotesize
\item Notes: All regressions contain the eight position indicators and firm
controls. The first four specifications include year fixed effects. The
last includes event fixed effects.
Stars use event-level CR2 $p$-values:
\sym{*} $p<0.10$, \sym{**} $p<0.05$, and \sym{***} $p<0.01$.
\end{tablenotes}
\end{threeparttable}
\end{table}

Dropping one of the 20 hardware securities at a time produces coefficients
between 2.59 and 3.69\pp{}. All 20 estimates are positive and significant
under event-level CR2 inference. Dropping NVIDIA yields 3.32\pp{}
($p=0.0028$). Dropping Sandisk produces the smallest coefficient,
2.59\pp{} ($p=0.0117$).

\section{Pairwise Position Contrasts}

Table~\ref{tab:a3} tests equality between the hardware coefficient and the
other ecosystem positions. The hardware--competitor difference is close to
zero. The hardware coefficient exceeds the average downstream and average
nonhardware coefficients under both return benchmarks and both inference
methods.

\begin{table}[H]
\centering
\begin{threeparttable}
\caption{Pairwise contrasts with the hardware coefficient}
\label{tab:a3}
\scriptsize
\setlength{\tabcolsep}{2.8pt}
\begin{tabular}{lrrrrrr}
\toprule
& \multicolumn{3}{c}{Market model}
& \multicolumn{3}{c}{Three-factor model} \\
\cmidrule(lr){2-4}\cmidrule(lr){5-7}
Hardware minus & Difference & Event $p$ & Two-way $p$ &
Difference & Event $p$ & Two-way $p$ \\
\midrule
Cloud platform        & 2.426\sym{**} & 0.046 & 0.072 & 2.257\sym{*} & 0.063 & 0.103 \\
Integrator            & 2.160\sym{**} & 0.025 & 0.051 & 2.243\sym{**} & 0.020 & 0.048 \\
Deployer              & 2.214\sym{**} & 0.013 & 0.004 & 2.162\sym{**} & 0.014 & 0.007 \\
Enabler               & 3.224\sym{***} & 0.004 & 0.014 & 3.561\sym{***} & 0.002 & 0.008 \\
Competitor            & 0.090 & 0.917 & 0.924 & 0.057 & 0.946 & 0.957 \\
Downstream average    & 2.533\sym{***} & 0.004 & 0.009 & 2.655\sym{***} & 0.002 & 0.007 \\
Nonhardware average   & 2.023\sym{**} & 0.012 & 0.022 & 2.056\sym{***} & 0.010 & 0.019 \\
\bottomrule
\end{tabular}
\begin{tablenotes}\footnotesize
\item Notes: Differences are percentage points. The downstream average
combines cloud, integrator, deployer, and enabler coefficients. The
nonhardware average combines the other seven position coefficients.
Stars use event-level CR2 $p$-values:
\sym{*} $p<0.10$, \sym{**} $p<0.05$, and \sym{***} $p<0.01$.
\end{tablenotes}
\end{threeparttable}
\end{table}

Holm adjustment treats the eight position coefficients as one testing
family. The adjusted $p$-values are 0.0119 for hardware and 0.0021 for
competitor. The remaining six adjusted $p$-values exceed 0.41.

\section{Time Interactions}

The time specification includes each of the eight position indicators, all
eight interactions with a 2025--2026 indicator, firm controls, and year
fixed effects. The eight post-2025 interactions are jointly significant
under every return-and-inference combination. Market-model $p$-values are
0.0048 with event CR2 and 0.0148 with event-and-firm clustering.
Three-factor $p$-values are 0.0017 and 0.00002.

Table~\ref{tab:a4} reports the formal difference between the hardware and
competitor changes. The estimate is approximately 5\pp{} under both return
models. All confidence intervals exclude zero.

\begin{table}[H]
\centering
\begin{threeparttable}
\caption{Hardware change minus competitor change after 2025}
\label{tab:a4}
\small
\setlength{\tabcolsep}{4pt}
\begin{tabular}{llrrrr}
\toprule
Return model & Inference & Difference & SE & 95\% CI & $p$ \\
\midrule
Market model
 & Event CR2 & 4.863\sym{***} & 1.468 & [1.954,\ 7.773] & 0.0013 \\
 & Event $\times$ firm & 4.863\sym{**} & 2.039 & [0.755,\ 8.972] & 0.0214 \\
\addlinespace
Three-factor model
 & Event CR2 & 5.032\sym{***} & 1.519 & [2.022,\ 8.043] & 0.0013 \\
 & Event $\times$ firm & 5.032\sym{**} & 2.117 & [0.766,\ 9.299] & 0.0219 \\
\bottomrule
\end{tabular}
\begin{tablenotes}\footnotesize
\item Notes: Difference and standard error are in percentage points. The
contrast is
$\Delta\beta_{\mathrm{hardware}}-\Delta\beta_{\mathrm{competitor}}$,
where each $\Delta$ is the coefficient on the position's interaction with
the 2025--2026 indicator.
Stars use the $p$-value in each row:
\sym{*} $p<0.10$, \sym{**} $p<0.05$, and \sym{***} $p<0.01$.
\end{tablenotes}
\end{threeparttable}
\end{table}

\section{Capability Specifications}

Table~\ref{tab:a5} reports the capability regressions. The pooled slope is
negative under both return models. The market-model slope is more negative
among closed-source events. The capability--hardware interaction is positive
under both return models and both inference methods.

\begin{table}[H]
\centering
\begin{threeparttable}
\caption{Capability slopes and hardware interactions}\label{tab:a5}
\scriptsize
\setlength{\tabcolsep}{3.2pt}
\begin{tabular}{llrrrrr}
\toprule
Return model & Sample or term & Coef. & Event $p$ & Two-way $p$ &
Observations & Events \\
\midrule
Market model
 & Capability, all events        & $-$0.0588\sym{*} & 0.071 & 0.082 & 3,842 & 87 \\
 & Capability, closed source     & $-$0.1305\sym{***} & 0.005 & 0.017 & 2,207 & 50 \\
 & Capability, open weight       & $-$0.0670 & 0.247 & 0.215 & 1,635 & 37 \\
 & Capability, unrelated firms   & $-$0.0814 & 0.623 & 0.507 & 174 & 87 \\
 & Hardware at mean capability   & 1.8527\sym{**} & 0.046 & 0.068 & 3,842 & 87 \\
 & Capability $\times$ hardware  & 0.1536\sym{**} & 0.045 & 0.051 & 3,842 & 87 \\
\addlinespace
Three-factor model
 & Capability, all events        & $-$0.0802\sym{**} & 0.030 & 0.030 & 3,842 & 87 \\
 & Hardware at mean capability   & 1.9497\sym{**} & 0.038 & 0.056 & 3,842 & 87 \\
 & Capability $\times$ hardware  & 0.1623\sym{**} & 0.034 & 0.039 & 3,842 & 87 \\
\bottomrule
\end{tabular}
\begin{tablenotes}\footnotesize
\item Notes: Coefficients are percentage points per intelligence-index
point, except for the hardware indicator. Capability has a standard
deviation of 12.7 points. Interaction regressions center capability before
forming the product term.
Stars use event-level CR2 $p$-values:
\sym{*} $p<0.10$, \sym{**} $p<0.05$, and \sym{***} $p<0.01$.
\end{tablenotes}
\end{threeparttable}
\end{table}

\section{Competitor Composition Tests}

Competitor observations cover AAPL, ADBE, GOOG, GOOGL, META, and MSFT.
Table~\ref{tab:a6} checks share classes, overlapping labels, and fixed
effects. Combining GOOG and GOOGL into one issuer cluster leaves the
two-way $p$-value at 0.022. Dropping either share class or the entire
Alphabet issuer produces the same coefficient.

\begin{table}[H]
\centering
\begin{threeparttable}
\caption{Competitor composition and specification checks}\label{tab:a6}
\scriptsize
\setlength{\tabcolsep}{3.4pt}
\begin{tabular}{lrrr}
\toprule
Specification & Coef. (pp) & Event CR2 $p$ & Two-way $p$ \\
\midrule
Baseline market model                    & 3.385\sym{***} & 0.0003 & 0.022 \\
Cluster GOOG and GOOGL as one issuer     & 3.385\sym{***} & 0.0003 & 0.022 \\
Drop GOOG                                & 3.396\sym{***} & 0.0003 & 0.022 \\
Drop GOOGL                               & 3.396\sym{***} & 0.0003 & 0.022 \\
Drop Alphabet issuer                     & 3.468\sym{***} & 0.0003 & 0.023 \\
Exclude competitor--cloud overlaps       & 3.478\sym{***} & 0.0003 & 0.023 \\
Exclude competitor--integrator overlaps  & 4.632\sym{***} & $<0.001$ & 0.007 \\
Exclude competitor--deployer overlaps    & 3.005\sym{***} & 0.007 & 0.074 \\
Event and firm fixed effects             & 1.320 & 0.414 & 0.032 \\
\bottomrule
\end{tabular}
\begin{tablenotes}\footnotesize
\item Notes: The final specification absorbs time-invariant firm-level
positions and has a rank-deficient design matrix for four coefficients.
The competitor coefficient remains estimable.
Stars use event-level CR2 $p$-values:
\sym{*} $p<0.10$, \sym{**} $p<0.05$, and \sym{***} $p<0.01$.
\end{tablenotes}
\end{threeparttable}
\end{table}

In six leave-one-competitor-security regressions, coefficients range from
3.00 to 4.28\pp{}. All six are significant under event CR2; five are
significant under two-way clustering. Dropping AAPL produces the largest
two-way $p$-value, 0.074. Issuer-level deletion gives the same pattern.

Deleting one event creator at a time produces 25 additional regressions.
Competitor coefficients range from 3.11 to 3.85\pp{}. Every estimate is
significant under event CR2 and two-way clustering. The two-way $p$-values
range from 0.011 to 0.034.

\section{Dependence and Overlapping Events}

Table~\ref{tab:a7} applies calendar-HAC adjustments to release-level
coefficients. Each release receives equal weight. Hardware estimates retain
their sign and magnitude as the bandwidth increases. Competitor estimates
also remain positive, with wider confidence intervals at longer
bandwidths.

\begin{table}[H]
\centering
\begin{threeparttable}
\caption{Event-level calendar-HAC inference}\label{tab:a7}
\small
\setlength{\tabcolsep}{4pt}
\begin{tabular}{lrrrrrr}
\toprule
Position & Coef. (pp) & Lag 0 & Lag 5 & Lag 10 & Lag 20 & Lag 40 \\
\midrule
Hardware   & 3.283 & 0.002 & 0.025 & 0.044 & 0.071 & 0.102 \\
Competitor & 2.118 & 0.019 & 0.072 & 0.100 & 0.124 & 0.142 \\
\bottomrule
\end{tabular}
\begin{tablenotes}\footnotesize
\item Notes: Entries after the coefficient are two-sided $p$-values.
Bandwidths are trading days.
\end{tablenotes}
\end{threeparttable}
\end{table}

A maximal set of non-overlapping twenty-day windows contains 29 events.
Across 250 random maximal sets, the median hardware coefficient is
4.22\pp{} and the median competitor coefficient is 3.78\pp{}. Every
draw produces positive coefficients for both positions. With event-and-firm
clustering, 49.6\% of hardware draws and 64.0\% of competitor draws have
$p<0.05$; the corresponding shares with $p<0.10$ are 82.0\% and 86.0\%.

Calendar-time portfolios give each trading day one observation. The
hardware post-minus-pre daily alpha is 7.50 basis points, with Newey--West
$p$-values of 0.039, 0.039, 0.056, and 0.092 at lags 5, 10, 20, and 40.
The competitor post-minus-pre alpha is 4.66 basis points, with $p$-values
from 0.218 to 0.308. These portfolio results complement the firm-event
regressions by changing both event and calendar-day weights.

\section{Abnormal Trading Volume}

Table~\ref{tab:a8} reports abnormal trading volume around model releases.
The dependent variable is log volume relative to its
$[-200,-11]$ estimation-window mean, scaled by the estimation-window
standard deviation. Hardware and competitor announcement-volume
coefficients are statistically insignificant in the NDXT45 sample.

\begin{table}[H]
\centering
\begin{threeparttable}
\caption{Abnormal volume before and after releases}\label{tab:a8}
\scriptsize
\setlength{\tabcolsep}{2.8pt}
\begin{tabular}{lrrrrrr}
\toprule
& \multicolumn{3}{c}{Pre-event $[-10,-2]$}
& \multicolumn{3}{c}{Announcement $[0,+1]$} \\
\cmidrule(lr){2-4}\cmidrule(lr){5-7}
Position & Coef. & Event $p$ & Two-way $p$ &
Coef. & Event $p$ & Two-way $p$ \\
\midrule
Hardware supplier      & $-$0.077\sym{*} & 0.051 & 0.353 & $-$0.007 & 0.882 & 0.943 \\
Cloud platform         & 0.074 & 0.148 & 0.299 & 0.095 & 0.161 & 0.400 \\
Downstream integrator  & $-$0.041 & 0.291 & 0.584 & $-$0.052 & 0.340 & 0.605 \\
Downstream deployer    & $-$0.112\sym{***} & 0.008 & 0.075 & $-$0.170\sym{***} & 0.004 & 0.054 \\
Downstream enabler     & $-$0.333\sym{***} & $<0.001$ & $<0.001$ & $-$0.457\sym{***} & $<0.001$ & $<0.001$ \\
Competitor             & 0.012 & 0.792 & 0.914 & 0.103 & 0.108 & 0.432 \\
Investor               & $-$0.097 & 0.232 & 0.436 & $-$0.063 & 0.603 & 0.757 \\
Listed owner           & $-$0.058 & 0.646 & 0.690 & 0.563\sym{***} & 0.007 & 0.032 \\
\bottomrule
\end{tabular}
\begin{tablenotes}\footnotesize
\item Notes: Regressions include all position indicators, firm controls,
and event fixed effects. Coefficients are in estimation-window standard
deviations.
Stars use event-level CR2 $p$-values:
\sym{*} $p<0.10$, \sym{**} $p<0.05$, and \sym{***} $p<0.01$.
\end{tablenotes}
\end{threeparttable}
\end{table}

\section{Relationship Codebook}

The coding unit is a firm--model-developer pair. Event indicators take the
union across developers represented in the event. Multiple indicators can
equal one for the same pair.

\paragraph{Hardware supplier.}
The firm earns material revenue from physical inputs used in model training
or inference, including accelerators, CPUs, memory, networking equipment,
servers, storage, semiconductor design tools, and fabrication equipment.

\paragraph{Cloud platform.}
The firm operates hyperscale computing infrastructure that rents training
or inference capacity and hosts third-party foundation models.

\paragraph{Downstream integrator.}
The firm embeds third-party foundation models into its own software product
or platform. The model is an input into a product sold under the listed
firm's brand.

\paragraph{Downstream deployer.}
The firm uses generative AI inside a broader operating business. The
firm's primary product is outside the foundation-model and AI-infrastructure
markets.

\paragraph{Downstream enabler.}
The firm advises, implements, integrates, or operates AI systems for
enterprise clients as a service.

\paragraph{Competitor.}
The firm develops foundation models in the same modality as the releasing
developer. Modality matching distinguishes language, image, video, audio,
and multimodal competition.

\paragraph{Investor and listed owner.}
The investor indicator records a financial stake in the releasing entity.
The owner indicator marks a listed parent or publisher responsible for the
released model.

\section{Reproducibility Map}

The canonical event-firm panel is
\path{Analysis/processed/event_firm_panel.csv}. Main position and
capability estimates are generated by
\path{Analysis/scripts/frl_two_way_cluster_and_equivalence.R}.
Pairwise contrasts are generated by
\path{Analysis/scripts/frl_position_contrasts.R}.
Time interactions and competitor checks are generated by
\path{Analysis/scripts/frl_time_and_competitor_robustness.R}.
Non-overlapping samples and calendar-time portfolios are generated by
\path{Analysis/scripts/frl_nonoverlap_calendar_portfolio.R}.
The root entry point \path{run_reproduction.sh} rebuilds the source-derived
data, the event-firm panel, every reported result, both coding validations,
the figure, and both PDFs.

\end{document}
